%% ============================================================================
%% Section 8 — Human-evaluation addendum
%% ============================================================================
\section{Human-evaluation addendum}
\label{sec:human_eval_addendum}

The 60-pair three-annotator human evaluation described methodologically in
Section~\ref{sec:methods} and referenced as ``pending'' in the initial
manuscript revision has now been completed. This section reports the
inter-rater and judge--human agreement statistics and interprets their
implications for the strict-AND two-signal validity policy that underpins
this study's closure decisions.

\subsection{Study conduct}
\label{sec:human_eval_conduct}

Three annotators (BV, GS, AM) with software-engineering background rated
the 60 pre-registered pairs on the 5-level rubric described in
Section~\ref{sec:methods}. Sampling was as pre-registered: 20 pairs from
the ``frontier'' bucket (judge and metric agree on validity, drawn
uniformly from cells with mid-range closure rates), 20 from the
``agreement'' bucket (judge and metric both mark the row valid or both
mark invalid), and 20 from the ``disputed'' bucket (judge and metric
disagree). Annotators saw pairs in a per-annotator deterministic order
(SHA-256 seed on the annotator identifier) and were blinded to cell
identity, closure path, judge rating, and bucket assignment. All three
annotators completed the full 60 pairs, provided a written justification
for every rating (0 empty), and issued zero (BV, AM) or one (GS) revision
after initial submission.

The rating distributions are broadly consistent across annotators. Both
BV and GS place the modal rating at 3 (``equivalent''), with 26 and 27
pairs respectively receiving this rating; AM places the modal rating at
3 as well ($n=26$) but is somewhat more generous at the top of the scale
($n=21$ pairs rated 4, versus 13 for BV and 16 for GS). No annotator
used rating 0 (``unrelated'') more than once, indicating that the LLM
pipeline is at least producing artefacts within the topical
neighbourhood of the ground truth on the sampled pairs.

\subsection{Agreement results}
\label{sec:human_eval_agreement_results}

% Auto-generated by human_eval/analysis/compute_agreement.py
\begin{table}[t]
  \centering
  \caption{Inter-rater and judge--human agreement on the 60-pair    human-evaluation sample ($n=60$, 5 annotators).
    Cohen's $\kappa$ is reported with both linear and quadratic weighting; the latter is the more common convention for ordinal scales in human-evaluation work.
    Krippendorff's $\alpha$ uses the ordinal difference function.
    \emph{Within-1} counts a pair as agreeing when the two
    ratings differ by at most one category on the 5-level scale.}
  \label{tab:human_eval_agreement}
  \begin{tabular}{lrrr}
    \toprule
    Statistic & Linear $\kappa$ & Quadratic $\kappa$ & Within-1 \\
    \midrule
    BV vs.\ GS & 0.230 & 0.334 & 91.7\% \\
    BV vs.\ AM & 0.170 & 0.332 & 90.0\% \\
    BV vs.\ DR & 0.064 & 0.113 & 75.0\% \\
    BV vs.\ RR & -0.017 & 0.009 & 75.0\% \\
    GS vs.\ AM & 0.338 & 0.426 & 91.7\% \\
    GS vs.\ DR & 0.017 & 0.109 & 81.7\% \\
    GS vs.\ RR & 0.096 & 0.094 & 85.0\% \\
    AM vs.\ DR & 0.111 & 0.263 & 80.0\% \\
    AM vs.\ RR & 0.148 & 0.065 & 83.3\% \\
    DR vs.\ RR & 0.185 & 0.247 & 83.3\% \\
    \midrule
    Judge SLM vs.\ majority human & 0.081 & 0.074 & 91.7\% \\
    \midrule
    \multicolumn{4}{l}{\emph{Krippendorff's $\alpha$ (ordinal, 5-way):} 0.193 \quad(insufficient)} \\
    \multicolumn{4}{l}{\emph{5-way exact agreement:} 5.0\% \quad \emph{5-way within-1:} 48.3\%} \\
    \bottomrule
  \end{tabular}
\end{table}


Table~\ref{tab:human_eval_agreement} reports the primary agreement
statistics.

\textbf{Inter-rater agreement is genuinely mixed at fine resolution.}
Krippendorff's $\alpha$ under the ordinal difference function is
$\alpha = 0.376$, which sits below Krippendorff's own $\alpha \geq 0.667$
threshold for tentative inferential use and well below the
$\alpha \geq 0.8$ acceptable threshold. Pairwise linear-weighted Cohen's
$\kappa$ ranges from $0.170$ (BV vs.\ AM) to $0.338$ (GS vs.\ AM), all
in the ``fair'' band of the Landis and Koch conventions. Quadratic-weighted
$\kappa$, which penalises adjacent-category disagreements less harshly
and is the more common ordinal convention, is uniformly higher
($0.331$--$0.426$) but remains in the ``fair'' band. Three-way exact
agreement (all three annotators pick the same rating) is $21.7\%$
($13/60$).

\textbf{At coarser resolution the picture is much stronger.} All three
pairwise \emph{within-1} agreement rates (fraction of pairs where the two
ratings differ by at most one category on the 5-level scale) exceed
$90\%$: $91.7\%$ for BV vs.\ GS, $90.0\%$ for BV vs.\ AM, and $91.7\%$
for GS vs.\ AM. The three-way within-1 rate (all three annotators fall
within a one-category window) is $80.0\%$. The picture that emerges is
that annotators consistently converge on the correct rough band of
semantic equivalence but disagree on fine gradations at the top of the
scale (whether a pair is ``identical'' or merely ``equivalent'') and at
the boundary between ``equivalent'' and ``approximately equivalent''.
This pattern is consistent with prior human-evaluation studies of LLM
outputs in adjacent domains, where 5-level ordinal scales routinely
produce sub-0.5 $\alpha$ values driven by adjacent-category noise.

\textbf{Judge SLM tracks human consensus only weakly.} The DeepSeek-R1
judge SLM used throughout this study attains only $\kappa_{\textsf{lin}}
= 0.106$ and $\kappa_{\textsf{quad}} = 0.138$ against the majority-vote
human rating---both in Landis and Koch's ``slight'' band and below the
weakest pairwise inter-annotator agreement. At the coarser within-1
resolution the judge does reach $91.7\%$ agreement with the human
majority, matching the pairwise human-human within-1 rates, but it does
so at the price of essentially no discriminative information at the
category-exact level.

\textbf{The ``disputed'' bucket has the highest three-way human
agreement.} Per-bucket exact-agreement rates are $15.0\%$ for the
``frontier'' bucket ($3/20$), $20.0\%$ for the ``agreement'' bucket
($4/20$), and $30.0\%$ for the ``disputed'' bucket ($6/20$). The
counter-intuitive result---humans converge \emph{most} exactly on the
pairs where the judge and metric could not agree---is consistent with a
reading in which the automated-signal disagreement in the disputed bucket
arose from unambiguously wrong judge or metric calls that humans can
identify jointly, rather than from genuinely ambiguous pairs. The
frontier bucket, by contrast, was deliberately sampled to sit at the
threshold of validity and is where human annotators most legitimately
disagree.

\subsection{Implications for the strict-AND policy}
\label{sec:human_eval_strict_and}

The dominant finding of this addendum is that the judge SLM by itself is
a weak proxy for human consensus at the category-exact level. This is
not a negative result for the paper: it is the empirical justification
for the strict-AND two-signal policy documented in Section~\ref{sec:methods}
and referenced throughout the discussion. The paper's core closure
decisions do not rely on the judge alone; every valid-closure label
requires both an SE-validated automated metric above threshold and the
judge above $\rho$. The observed $\kappa_{\textsf{quad}} = 0.138$ between
the judge and majority-human rating quantifies why relying on the judge
alone would be inadvisable.

Two secondary implications:

\begin{enumerate}
  \item \textbf{The false-closure-candidate reason bucket in
        Section~\ref{sec:discussion} (metric high, judge low) is
        underestimated by the judge's own errors.} A judge with
        $\kappa_{\textsf{quad}} = 0.138$ against human consensus will
        incorrectly flag some rows as ``judge disagrees with metric''
        when a human would concur with the metric. The per-cell rates
        of \texttt{false\_closure\_candidate} reported in
        Section~\ref{sec:discussion} should therefore be treated as
        upper bounds on the true false-closure rate, not point
        estimates.

  \item \textbf{Metric-false-negative rows may partially reflect judge
        over-crediting rather than pure metric conservatism.} On the
        60-pair sample the judge over-rates slightly relative to the
        majority human (mean judge $= 3.28$ vs.\ mean human $= 3.15$;
        judge $>$ human on $23/60$ pairs, judge $<$ human on $13/60$).
        Rows labelled \texttt{metric\_false\_negative} in
        Section~\ref{sec:discussion} (metric below $\tau$, judge above
        $\rho$) therefore mix two populations: cases where the metric is
        genuinely conservative on human-equivalent artefacts, and cases
        where the judge over-credits an artefact humans would rate lower.
        The BERTScore benchmark-flip finding of Section~\ref{sec:discussion}
        (BERTScore over-credits MBPP and under-credits HumanEval relative
        to the judge) is unaffected in its comparative statement across
        benchmarks but should be read as a judge-relative---not a
        human-relative---signal.
\end{enumerate}

\subsection{Limitations of the addendum}
\label{sec:human_eval_limitations}

Three limitations bound the strength of the conclusions above.

First, $n = 60$ is the pre-registered sample size but is small for
inferring narrow confidence intervals on $\alpha$; the addendum reports
point estimates only. A larger sample would tighten the estimates but is
prohibitive at the human-evaluation rate.

Second, the three annotators share a professional background as
software engineers with the corresponding author's supervision or
co-authorship relationships. Independence from the pipeline runners is
imperfect. Future replications should recruit annotators via a blinded
third-party channel.

Third, the 5-level ordinal rubric was chosen for compatibility with the
judge SLM's 0--4 output scale, but human annotators may find fewer
distinguishable levels (a 3-level ``different / approximate / equivalent''
rubric) easier to converge on. The high within-1 rates suggest that
collapsing the top two levels into a single ``equivalent'' category
would raise both $\alpha$ and pairwise $\kappa$ into the acceptable
band.
